%%%PAGE 017 rot=0 legibility=good summary=磁场中粒子圆轨迹技巧续：环形磁场、圆形磁场汇聚发散、类圆形磁场、双直线边界
%%%BLOCK id=017-01 ch=11 sec=环形磁场 kind=method alt=none cont=none y=0.06-0.29
\textbf{环形磁场}
%FIG p017_a 0.10 0.095 0.37 0.29
\notefig[0.3]{figures/p017_a.png}
最值相切\quad 径入径出
\[ (R-R_2)^2=R_2^2+r^2 \qquad R_2=\frac{mV_2}{qB}\ \Rightarrow\ V_2 \]
\[ (R_1+r)^2=R^2+R_1^2 \qquad R_1=\frac{mV_1}{qB}\ \Rightarrow\ V_1 \]
%%%END
%%%BLOCK id=017-02 ch=11 sec=圆形磁场·汇聚发散 kind=method alt=none cont=none y=0.29-0.60
\textbf{圆形磁场汇聚发散}\quad $R_{\text{轨迹}}=R_{\text{磁场}}$

点入平出\quad 平入点出
%FIG p017_b 0.095 0.368 0.352 0.595
\notefig[0.4]{figures/p017_b.png}
%FIG p017_c 0.41 0.362 0.67 0.545
\notefig[0.35]{figures/p017_c.png}
菱形 $\Rightarrow BO_2\parallel AO_1\Rightarrow V_{\text{末}}\parallel L$
%%%END
%%%BLOCK id=017-03 ch=11 sec=类圆形磁场 kind=method alt=none cont=none y=0.59-0.80
\textbf{类圆形磁场}（半圆、\unclear{$\frac14$圆}、扇形、树叶形）

粒子没有经过的磁场，是可有可无的磁场
%FIG p017_d 0.09 0.665 0.30 0.80
\notefig[0.25]{figures/p017_d.png}
入射范围角$=$圆心角
\[ S_{\mathrm{tot}}=2\left(S_{\text{扇}}-S_{\triangle}\right) \]
% CHECK: 下标手写为"ot"（可能是 tot 或 总）；S 的第二个下标像 角/扇，按物理意义取 扇
%%%END
%%%BLOCK id=017-04 ch=11 sec=双直线边界 kind=method alt=none cont=next y=0.80-1.0
\textbf{双直线边界}

最值相切\quad 能从上口射出
%FIG p017_e 0.04 0.884 0.208 1.0
\notefig[0.25]{figures/p017_e.png}
定向不定速

$\checkmark$最\unclear{小}（最\unclear{小}$V$）\quad 圆的放缩
% CHECK: 第一个对勾后的"最?"与括号内"最?V"手写潦草，按物理意义（相切对应最小速度）取 小
能从下口射出的

$\checkmark$最大（最\unclear{大}$V$）
% CHECK: 括号内"最?V"手写潦草，像 木（大/小 均有可能），按物理意义取 大
%%%END
%%%PAGEEND 017
